\section{Curve zeta functions in Hall operator exchange}
\label{intake:curve:section}

For a curve over a finite field, the local factors also occur as
coefficients in an operator exchange. The underlying algebra now comes
from extensions of torsion sheaves, rather than the divisor monoid.
A torsion sheaf on a smooth projective curve is the direct sum of its
parts supported at closed points. This decomposition turns local Hall
operators into global operators. Their exchange coefficient counts the
points over every finite extension of the ground field.

Let $X$ be a smooth projective geometrically connected curve over
$\mathbb F_q$. For a closed point $x$, put $d_x=[\mathbb F_q(x):\mathbb F_q]$
and $Q_x=q^{d_x}$. The completion of the local ring is
$\mathbb F_{Q_x}[[t]]$, after a choice of parameter and coefficient field.
Its finite-length modules are nilpotent modules over $\mathbb F_{Q_x}[t]$.

\subsection{Local Hall operators and their normalization}

For a prime power $Q$, let $H_Q$ be the complex Hall algebra of these
nilpotent modules. Its basis element $u_M$ denotes the isomorphism class
of $M$. The coefficient of $u_C$ in $u_Au_B$ counts submodules
$B'\subset C$ with $B'\simeq B$ and $C/B'\simeq A$.
Give it the positive Hermitian form, conjugate-linear in the first slot,
\[
 \langle u_M,u_N\rangle=\frac{\delta_{MN}}{|\operatorname{Aut}M|}.
\]
The Green coproduct is adjoint to multiplication for this form.
In this category it is cocommutative and multiplicative.
These are the classical Hall bialgebra identities, on this untwisted
category with zero Euler form. The polynomial presentation below is
Steinitz--Hall's theorem; see Schiffmann, \emph{Lectures on Hall algebras},
Theorem~2.6 and equation~(2.9).

Write $u_{1^r}$ for the zero endomorphism of $\mathbb F_Q^r$ and set
\[
 e_r=Q^{r(r-1)/2}u_{1^r},\qquad
 E(z)=\sum_{r\ge0}e_rz^r,\qquad e_0=1.
\]
Subspace counting gives
$\Delta u_{1^r}=\sum_{i+j=r}Q^{-ij}u_{1^i}\otimes u_{1^j}$.
Thus $\Delta E(z)=E(z)\otimes E(z)$. The elements $e_r$, $r>0$,
freely generate $H_Q$. One can see the triangular presentation directly:
for a Jordan partition $\lambda$, use the successive dimensions of
$\ker t^i/\ker t^{i-1}$. The product of the corresponding semisimple
classes counts filtrations lowered by $t$. Its coefficient on $u_\lambda$
is one, since equality forces the kernel filtration. Every other term
has larger kernel dimensions. The resulting triangular matrix is
invertible in each finite dimension.

Define $p_r$ by the formal identity
\[
 \log E(z)=\sum_{r\ge1}\frac{(-1)^{r-1}}r p_rz^r.
\]
They are primitive polynomial generators. Their norms can be computed
without a choice of an orthonormal basis. Since
$|\operatorname{GL}_r(\mathbb F_Q)|=Q^{r(r-1)/2}\prod_{j=1}^r(Q^j-1)$,
\[
 \langle E(u),E(v)\rangle
 =\sum_{r\ge0}\frac{Q^{r(r-1)/2}(\bar uv)^r}{\prod_{j=1}^r(Q^j-1)}
 =\prod_{j\ge1}(1-Q^{-j}\bar uv)^{-1}.
\]
The last equality follows by comparing coefficients in
$ (1-t)P(Qt)=P(t)$, with constant coefficient one.
Adjunction and primitivity make distinct primitive degrees orthogonal
and give factorial pairings for their powers. Taking logarithms therefore gives
\[
 \langle p_r,p_s\rangle=\delta_{rs}\frac{r}{Q^r-1}.
\]
Consequently the following operators are adjoint on the polynomial domain:
\[
 a_{-r}=\sqrt{Q^r-1}\,M_{p_r},\qquad
 a_r=\frac{r}{\sqrt{Q^r-1}}\frac{\partial}{\partial p_r},\qquad
 [a_r,a_{-s}]=r\delta_{rs}I.
\]
Here $M_{p_r}$ means multiplication by $p_r$, and the square root is positive.
No assertion that these unbounded operators act everywhere on the Hilbert
completion is needed.

\subsection{The global exchange}

The torsion Hall algebra of $X$ is the algebraic tensor product
$H_X^{\mathrm{tor}}=\bigotimes_{x\in|X|}^{\mathrm{alg}}H_{Q_x}$,
with vacuum $1$ in all but finitely many factors. Distinct supports have
no extensions or homomorphisms between them. Their Hall products and
pairings therefore factor. Give a length-$r$ local module degree $rd_x$.
There are finitely many closed points of bounded degree, so each
homogeneous piece of this algebra is finite-dimensional.
Let $a_{x,\pm r}$ be the preceding operators on the $x$ factor, and define
\[
 A_{-n}=\sum_{d\mid n}d\sum_{d_x=d}a_{x,-n/d},\qquad
 A_n=A_{-n}^{\dagger}\qquad(n>0).
\]
These are finite sums on the common polynomial domain. A closed point
of degree $d$ contributes $d$ geometric points over $\mathbb F_{q^n}$
exactly when $d\mid n$. Thus, writing $N_n=\#X(\mathbb F_{q^n})$,
\[
 [A_n,A_{-m}]
 =\delta_{nm}\sum_{d\mid n}d^2\frac nd\#\{x:d_x=d\}I
 =\delta_{nm}nN_nI.
\]
All commutators between two positive or two negative modes vanish.

Put $s=w^{-1}$, with $z,s$ independent formal variables, and set
\[
 \Gamma_-(z)=\exp\left(\sum_{n>0}\frac{A_{-n}}n z^n\right),\qquad
 \Gamma_+(w)=\exp\left(\sum_{n>0}\frac{A_n}n w^{-n}\right).
\]
Each coefficient is a finite sum of polynomial-domain operators.
Products are identities in
$\operatorname{End}_{\mathbb C}(H_X^{\mathrm{tor}})[[z,s]]$.
For the formal zeta function
\[
 Z_X(t)=\exp\left(\sum_{n>0}\frac{N_n}n t^n\right)
       =\prod_{x\in|X|}(1-t^{d_x})^{-1},
\]
the logarithm of the Euler product proves the second equality coefficientwise.
The central commutator of the two exponents equals $\log Z_X(zs)$.
Expanding the exponentials, or differentiating their formal conjugation
identity, therefore proves
\[
 \Gamma_+(w)\Gamma_-(z)
 =Z_X(z/w)\Gamma_-(z)\Gamma_+(w).
\]
This is an operator identity before taking a trace. Repeated reordering
uses the product of the central factors for all exchanged pairs.
The statement specifies a formal expansion at $z/w=0$; an analytic
operator identity would require additional domain and convergence estimates.

For $X=\mathbb P^1$, counting gives $N_n=q^n+1$, hence
$Z_X(t)=((1-t)(1-qt))^{-1}$.
For the smooth projective cubic
$E/\mathbb F_2: y^2+y=x^3+x+1$, the first four point counts are
$1,5,13,25$. The trace criterion for $y^2+y=b$ gives two solutions
when $\operatorname{Tr}_{\mathbb F_{2^n}/\mathbb F_2}(b)=0$, and none otherwise.
For $n=1,2,3,4$, it gives respectively $0,2,6,12$ admissible $x$ values,
with one additional point at infinity. Thus
\[
 Z_E(t)=1+t+3t^2+7t^3+15t^4+O(t^5).
\]
These coefficients determine the corresponding operator exchange to degree four.
The full rational formula
$Z_E(t)=(1-2t+2t^2)/((1-t)(1-2t))$
uses the classical elliptic-curve zeta theorem in addition to these finite counts.

\subsection{Torsion acting on a line bundle}

The full Hall algebra of coherent sheaves also contains the classes of
line bundles. Fix a rational point $x$ and a line bundle $L$.
There are two middle sheaves in an extension of $\mathcal O_x$ by $L$:
the split sheaf $L\oplus\mathcal O_x$, and $L(x)$.
Indeed $\operatorname{Ext}^1(\mathcal O_x,L)$ is one-dimensional,
and all nonzero classes have the latter middle sheaf.
In the split middle sheaf, the relevant copies of $L$ are the graphs of
the $q$ maps $L\to\mathcal O_x$. In $L(x)$ the relevant subobject is
uniquely $L(x)(-x)=L$. Conversely, an extension of $L$ by $\mathcal O_x$
splits because $\operatorname{Ext}^1(L,\mathcal O_x)=0$.
Its unique torsion subobject is $\mathcal O_x$. Therefore
\[
 u_{\mathcal O_x}u_L=q\,u_{L\oplus\mathcal O_x}+u_{L(x)},\qquad
 u_Lu_{\mathcal O_x}=u_{L\oplus\mathcal O_x}.
\]
Let $v=\sqrt q$ and define the twisted product
$u_A*u_B=v^{\chi(A,B)}u_Au_B$, where
$\chi(A,B)=\dim\operatorname{Hom}(A,B)-\dim\operatorname{Ext}^1(A,B)$.
Here $\chi(\mathcal O_x,L)=-1$ and $\chi(L,\mathcal O_x)=1$.
The ordinary commutator for this twisted product consequently satisfies
\[
 [u_{\mathcal O_x},u_L]_* =v^{-1}u_{L(x)}.
\]
On $\mathbb P^1$, all rational-point modifications of $\mathcal O(n)$
are isomorphic to $\mathcal O(n+1)$. Summing over the $q+1$ rational points gives
\[
 \left[\sum_{x\in\mathbb P^1(\mathbb F_q)}u_{\mathcal O_x},
 u_{\mathcal O(n)}\right]_*
 =(v+v^{-1})u_{\mathcal O(n+1)}.
\]
This mixed Hall operation retains the elementary modification of the line
bundle. Identifying a full quantum loop algebra additionally requires
its remaining generators and relations. Neither this mixed operation nor
the torsion exchange constructs a comparison with a global arithmetic
trace category or with the full Weil functional.
